




\subsection{Toom's Contours}
\label{subsection:toricencoded}
We are about to zoom out from a single register encoded in the Toric code to describe a fully fault-tolerant computation. For a given quantum circuit, assumed to consist of gates $\mathbf{X}, \mathbf{Z}, \mathbf{CX}, \mathbf{H}, \mathbf{S}, \mathbf{T}$ and measurements, we define its Toric encoded circuit as its version where each qubit is encoded by the Toric code. Formally:
    \begin{definition}[Toric Encoded Circuit]
      Let $C$ be a quantum circuit. We say that $C$ is a \textit{Toric encoded circuit} if the following holds:
      \begin{enumerate}
        \item Any layer of the circuit consists of classical and quantum registers, and there is a Toric code $\mathcal{C}$ such that the quantum registers are encoded by the \textit{nice encoding of} $\mathcal{C}$. At time $t=0$, the quantum registers are initialized to be either $\ket{\bar{0}}$, encoded $\ket{\mathbf{H}}$, encoded $\ket{\mathbf{S}}$, or encoded $\ket{\mathbf{T}}$. The classical registers are initialized to be zeros.
        \item The layers of $C$ on the quantum registers at even times are sweep-cycles, and the layers at odd times are realizations of either logical Paulis, $\mathbf{CX}$, or measurements in the computational and parity bases. On the classical registers, including the measured register, we allow any two-bit gates.
        \end{enumerate} 
    \end{definition}

For Toric encoded circuits, we define their \textit{base graph}, which can be thought of as a recording of a running from the viewpoint of the Toric complex. The \textit{base graph} is obtained by taking the Toric complex's vertices at different times as the vertices and connecting any two sequential in time if they are support faces — qubits — which are either the same or are entangled by a quantum gate.

  \begin{definition}[Base graph $G_{o}$]
    Let $C$ be a Toric encoded circuit at depth $\mu$. Notice that measurements take quantum registers into classical data; thus each quantum register has a final time coordinate, which we denote by $\{\mu_i\}_i$. Denote by $\{\mathcal{G}_{i}\}_i$ the skeletons of the $\infty$-lift of the Toric complexes that induce the codes. We define the base graph of $C$, denoted by $G_{o} = (V_{o}, E_{o})$, as follows:
    \begin{enumerate}
      \item The vertices $V_{o}$ are pairs, where the first coordinates correspond to a vertex in the disjoint union over the vertices in $\{\mathcal{G}_{i}\}$, and the second coordinate is associated with a time. Namely:
        \begin{equation*}
          \begin{split}
            V_{o} \colonequals \bigcup_{i} \bigl( V(\mathcal{G}_{i}) \times [\mu_{i}] \bigr)
          \end{split}
        \end{equation*}
      \item The edges $E_{o}$ are derived from the original edges in $\mathcal{G}_{1}, \ldots, \mathcal{G}_{k}$, with the addition that any two vertices in preceding times $t$ and $t+1$ that support qubits (faces) such that $C^{\infty}$ acts non-trivially on those qubits at time $t$ are also connected by an edge. Namely:
        \begin{equation*}
          \begin{split}
            E_{o} & \colonequals \big\{ \{(v,t), (u,t^{\prime}) \} :  \text{ if }     \\ 
              & \text{either there exists } i \text{ such }  \{v,u\} \in \mathcal{G}_{i} \text{ and } t^{\prime} = t+1 \\
            & \text{or there exists } u_i = (g_{i},l_{i}) \in C^{\infty} \text{ such }  v,u \in l_{i,2} \text{ and } t = l_{i,1}, t^{\prime} = l_{i,1}+1  \big\} 
          \end{split}
        \end{equation*}
Where we define a vertex $v$ to belong to location $l$ if there is a qubit-face $f$ in $l$ such that $v$ belongs to its associated face.
    \end{enumerate} 
  \end{definition}
Notice that since the gates' locations are invariant to the conditioning of outcomes measurement (\Cref{remark:condmeas}), we have that for any outcomes vector $\alpha \in \mathbb{F}_{2}^{(\cdot)}$, the base graph of the conditioned circuit $C[\mathcal{E},\alpha]$ is the same as for $C[\mathcal{E}]$.

\paragraph{}

As a concept, if the base graph is a recording of a run from the Toric complexes' point of view, then a \textit{Toom's contour} is a recording of the deviation's evolution in the eyes of the Toric complexes' vertices. We use the notion of $x,z$-component of a lifted Pauli and the Pauli-path to describe a deviation:
\begin{definition}
Let $P \in \mathcal{P}^{(\cdot)}$ be a Pauli, and let $Q \in \mathcal{P}^{(\cdot)}$ be its $\infty$-lift. Assume that $P$ acts on a Toric encoded circuit, namely it has a decomposition such that each element of it acts on a register encoded in the $(d,d^{\prime})$-Toric code. We denote by $x(Q)$ the $d^{\prime}$-dimensional chain in the infinite complex that supports the $X$-component of $Q$, and by $z(Q)$ the $d-d^{\prime}$-dimensional chain, such that the inverse $\phi$-transpose sends it to the $d^{\prime}$-dimensional chain that supports the $Z$-component of $Q$, namely $\phi^{-1}(z(Q))$ is the chain that supports the $Z$-component.



\end{definition}

\begin{definition}
  \label{def:paulipath}
  Consider the propagation $(\mathbf{P},\mathbf{a})$ of a conditioned subjected circuit $C[\mathcal{E},\alpha]$ (\Cref{def:propagation}). A Pauli-path of $C[\mathcal{E},\alpha]$ is the $\infty$-lift of the Paulis in $\mathbf{P}$.
 

\end{definition}


For a Pauli-path $P_{1},..,P_{\mu+1}$, the Toom's contour is the restriction of the base graph to vertices that witness a deviation, namely, those that belong to qubits supporting $P_{1},P_{2},..,P_{\mu+1}$. 


\begin{definition}[Toom's Contour Graph (Standard Form)]
    \label{def:contour}
    For a Toric encoded quantum circuit $C$, set of faults $\mathcal{E}$, measurement outcomes $\alpha$, and a Pauli-path $P_{1},..$ of $C[\mathcal{E}, \alpha]$. We define the \textit{Toom's contour graph} $\mathcal{H} = (V_{\mathcal{H}}, E_{\mathcal{H}})$ as follows: 
    \begin{enumerate}
      \item The vertices $V_{\mathcal{H}} = V_{\mathcal{H}, X} \cup V_{\mathcal{H}, Z}$ correspond to vertices of faces supporting the Pauli-path. They consist of the triple $(v,t,\cdot)$ where $(v,t)$ is a vertex of the base graph, and the third is a symbolic coordinate that is either $\{X\}$ or $\{Z\}$ corresponding to the component of the Pauli that $v$ supports at time $t$.
        \begin{equation*}
          \begin{split}
            V_{\mathcal{H}, X} \colonequals & \big\{ v_{t} \times \{X\} : v_{t}=(v,t) \in V_{o} \\
            & \, \, \, \text{ and there exists a face } f \in \partial^{-} x(P_{t}) \text{ such } v \in f \big\} \\
             V_{\mathcal{H}, Z} \colonequals & \big\{ v_{t} \times \{Z\} : v_{t}=(v,t) \in V_{o} \\
             & \, \, \, \text{ and there exists a face } f \in \partial^{-} z(P_{t}) \text{ such } v \in f \big\} 
          \end{split}
        \end{equation*}


      \item The edges $E_{\mathcal{H}}$ is partitioned into two types $A_{\mathcal{H}}$ and $F_{\mathcal{H}}$. For any realization of logical gate $g$ which applied in time $t$ in $C[\mathcal{E}]^{\infty} $, we denote by $E_{o}|_{g}$ the edges in $E_{o}$ which induced by $g$. Then, we add edge to $A_{\mathcal{H}}$ as follows: 

        \begin{enumerate}
          \item If $g$ does not belong to the faults, namely $g \notin \mathcal{E}$, and $g$ is not a measurement in the phases-basis, then we add the arrows: 
            \begin{equation*}
              \begin{split}
                 \bigl\{  \{ \tilde{u} = (u,t, \cdot),  \tilde{v} = (v,t+1, \cdot)\} : \{\{u,t\},\{v,t+1\}\} \in E_{o}|_{g} \text{ and } \tilde{v},\tilde{u} \in V_{\mathcal{H}, \cdot} \bigr\}
              \end{split}
            \end{equation*}
            Where if $g$ is a measurement in the phase-basis, then we add type-alternating arrows as follows:
            \begin{equation*}
              \begin{split}
                & \bigl\{  \{ \tilde{u} = (u,t, \{Z\}),  \tilde{v} = (v,t+1, \{X\})\} : \{\{u,t\},\{v,t+1\}\} \in E_{o}|_{g} \\
              &  \, \, \, \text{ and } \tilde{v}\in V_{\mathcal{H},X},\tilde{u} \in V_{\mathcal{H},Z} \bigr\} 
              \end{split}
            \end{equation*}
          \item Else, if $g \in \mathcal{E}$ then, we add:   
            \begin{equation*}
              \begin{split}
                & \bigl\{  \{ \tilde{u} = (v,t, \cdot),  \tilde{v} = (v,t+1, \cdot)\} : \text{ such either } \tilde{u} \in V_{\mathcal{H},X} \text { or } \tilde{u} \in V_{\mathcal{H},Z} \bigr\} 
              \end{split}
            \end{equation*}
        \end{enumerate}
            To construct $F_{\mathcal{H}}$, we connect any pair of vertices $(v,t)$, $(u,t)$ if there exists a vertex $(w,t+1)$ such both of them are connected to it through $A_{\mathcal{H}}$. Namely: 
        \begin{equation*}
          \begin{split}
            F_{\mathcal{H}} &= \big\{ e = \{u,v\} : \text{ exists } w \in V_{\mathcal{H}} \text{ such } \{v,w\},\{u,w\} \in A_{\mathcal{H}} \big\} 
          \end{split}
        \end{equation*}
    \end{enumerate}
    We call $A_{\mathcal{H}}$ and $F_{\mathcal{H}}$ the \textit{arrows} and \textit{forks} edges respectively. 
\end{definition}












For the analyses, we would like to extend Toom's contour to record in their final layer a partial time-step, namely a computation up to a certain step. For this, we first consider Toom's contour of the partial computation circuit up to a step (\Cref{def:steps}). Then, since we would like to refer to the connected chain at the final layer as the 'final error clusters', we add at the end a layer whose purpose is to connect in time (by arrows) the connected components of Toom's contour in the prior layer, which are connected by checks (their union is identified with a connected chain).

\begin{definition}[Analysis Toom's Contour Up to Step $\tau$]
  For a Toric encoded quantum circuit $C$, a set of faults $\mathcal{E}$, outcomes $\alpha$, and a step $\tau$, let $C[\mathcal{E},\alpha]_{\tau}$ be the partial computation circuit of $C[\mathcal{E},\alpha]$ up to step $\tau$. We define the \textit{analysis Toom's contour at step $\tau$} $\mathcal{H}_{\tau} = (V_{\mathcal{H}}, E_{\mathcal{H}})$ to be the graph obtained by taking the Toom's contour of $C[\mathcal{E},\alpha]_{\tau}$ and adding vertices and edges as follows:
  \begin{enumerate}
    \item Let $t$ be the depth of $C[\mathcal{E},\alpha]_{\tau}$. We add a layer corresponding to the time $t+1$, which duplicates the vertices of time $t$. Namely, we add to $V_{\mathcal{H}}$:
\begin{equation*}
  \begin{split}
    \big\{ (v,t+1, \cdot ) : \text{ if }  (v,t, \cdot ) \in V_{\mathcal{H}}  \big\}
  \end{split}
\end{equation*}
    \item Then, we add the following additional arrows:
        \begin{equation*}
          \begin{split}
            \bigl\{ \bigl\{ \tilde{u} & = (u,t, \cdot ),  \tilde{v}= (v,t+1,\cdot)  \bigr\}  \text{ if  either there exists }\\ 
           & \text{ a } d^{\prime}-1  \text{ dimensional face } f  \text{such } \tilde{u},\tilde{v} \in V_{\mathcal{H},X} \text{ and } u,v \in f  \\
          & \text{or there exists a } d - d^{\prime}-1  \text{ dimensional face } f \text{ such } \tilde{u},\tilde{v} \in V_{\mathcal{H},Z} \text{ and } v,u \in f   \bigr\}  
          \end{split}
        \end{equation*}
      \item Finally, we add forks to preserve the property that any two vertices with the same time coordinate, which are connected to a third (with a $+1$ higher time coordinate) via arrows, are also connected by a fork.
  \end{enumerate}
\end{definition}




